\documentclass[12pt]{article}
\usepackage{amsmath}
\usepackage{geometry}
\geometry{letterpaper, margin=1in}
\usepackage{fancyhdr}
\usepackage{setspace}
\usepackage{hyperref}

\pagestyle{fancy}
\fancyhf{}
\rhead{33-120 Science and Science Fiction}
\lhead{Sci-Fi Problem 3 (25 points)}
\setlength{\headheight}{15pt}

\begin{document}

\vspace{1cm}
\noindent \textbf{Film reference}: \textit{Despicable Me} (Chris Renaud and Pierre Coffin, Universal 2009)

\noindent \textbf{Concept}: an object of sufficient mass or sufficient density will have a so-called event horizon – a distance of approach, beyond which nothing can escape, not even light. This can be described mathematically by the Schwarzschild radius:

\begin{equation}
r_s = \frac{2Gm}{c^2}
\end{equation}

\noindent where $G = 6.67 \times 10^{-11} \, \text{Nm}^2/\text{kg}^2$ is the Universal gravitational constant, $m$ is the mass of the object, and $c$ is the speed of light. A black hole is an object whose physical size is smaller than its Schwarzschild radius. If the Schwarzschild radius lies somewhere inside the physical object, then there is no event horizon, and the object cannot be a black hole.

\vspace{0.5cm}

\noindent You will receive 5 points for an on time submission. The final 20
points are assigned below. In order to get full credit for the problems below,
be sure to:
\begin{itemize}
    \item explicitly state all assumptions made that are not given in the
    problem.
    \vspace{-10pt}
    \item show all of the steps in each calculation.
    \vspace{-10pt}
    \item include proper physical units for all results.
\end{itemize}

\vspace{10pt}
\noindent \textbf{Problem 1. (20 points)} \begin{enumerate}
    \item[a)] (4 pts) Go to your favorite source of reliable information and look up a value for the mass of the Moon. (Express the value in units of kilograms.)
    
    \item[b)] (8 pts) Calculate the Schwarzschild radius, \( r_s \), for the Moon, using the equation given above. (Express the result in units of meters.)
    
    \item[c)] (8pts) In the 2009 movie \textit{Despicable Me}, the arch-criminal Gru uses a “shrink ray” to reduce the Moon to roughly the size of a grapefruit or a softball. Is this shrunken size sufficiently small (assuming that the mass of the Moon has not been changed by shrinking it) for the Moon to become a black hole? Explain your answer.
\end{enumerate}

\newpage
\vspace{10pt}
\vspace{10pt}
\noindent \textbf{Problem 1. (20 points)} \begin{enumerate}
    \item[a)] (\textbf{4 pts}) Go to your favorite source of reliable information and look up a value for the mass of the Moon. (Express the value in units of kilograms.) \newline \newline $M_{\rm{Moon}}=7.35\times10^{22}\,\rm{kg}$. 
    
    \item[b)] (\textbf{10 pts}) Calculate the Schwarzschild radius, \( r_s \), for the Moon, using the equation given above. (Express the result in units of meters.) \newline \newline $R_{\rm{Schwarzschild,\,Moon}} = \frac{2GM_{\rm{Moon}}}{c^2} = \frac{2\times6.67\times10^{-11}\,\rm{m^3/(kg\,s^2)}\times7.35\times10^{22}\,\rm{kg}}{(3\times10^8\,\rm{m/s}^2)} =0.00011\,\rm{m} = 0.011\,\rm{cm}$
    
    \item[c)] (\textbf{6 pts}) In the 2009 movie \textit{Despicable Me}, the arch-criminal Gru uses a “shrink ray” to reduce the Moon to roughly the size of a grapefruit or a softball. Is this shrunken size sufficiently small (assuming that the mass of the Moon has not been changed by shrinking it) for the Moon to become a black hole? Explain your answer. \newline \newline No, a grapefruit is much larger than a centimeter across. 
\end{enumerate}

\end{document}
